%-------------------------
% CV of Jiaqi Xue
% Layout modeled on an academic CV: Palatino body, blue section headings,
% plain-text publication entries (own name in bold, no italics/underlines).
%-------------------------

\documentclass[letterpaper,10pt]{article}

\usepackage[T1]{fontenc}
\usepackage{mathpazo}
\usepackage[letterpaper,margin=0.75in]{geometry}
\usepackage{xcolor}
\usepackage{titlesec}
\usepackage{needspace}
\usepackage{enumitem}
\usepackage[hidelinks]{hyperref}
\input{glyphtounicode}
\pdfgentounicode=1   % keep the PDF machine readable / ATS parsable

\definecolor{secblue}{RGB}{31,105,160}
\definecolor{rulegray}{RGB}{150,150,150}

\pagestyle{empty}
\setlength{\parindent}{0pt}
\setlength{\parskip}{0pt}
\raggedbottom
\emergencystretch=1.5em   % avoid overfull lines around unbreakable names

% Section headings: blue title with a thin gray rule underneath
% Keep a heading with at least a few lines of its section (no orphaned headings)
\titleformat{\section}{\needspace{5\baselineskip}\color{secblue}\large\bfseries}{}{0em}{}[{\color{rulegray}\vspace{-1pt}\titlerule}]
\titlespacing*{\section}{0pt}{10pt}{4pt}

%-------------------------
% Custom commands

% Education / experience entry: #1 institution, #2 location, #3 role, #4 dates
\newcommand{\entry}[4]{%
  \textbf{#1}\hfill #2\\
  \textit{#3}\hfill #4\par\vspace{3pt}}

% One-line description under an entry
\newcommand{\entrydesc}[1]{\vspace{-2pt}#1\par\vspace{4pt}}

% Numbered publication with hanging indent: #1 number, #2 text
\newcommand{\pub}[2]{%
  \par\needspace{3\baselineskip}% keep short entries on one page
  \hangindent=2.1em\hangafter=1
  [#1]~#2\par\vspace{2.5pt}}

% Paper title linked to its PDF (links are not colored)
\newcommand{\ptitle}[2]{\href{#1}{#2}}

\newcommand{\me}{\textbf{Jiaqi Xue}}

%-------------------------------------------
\begin{document}

%----------HEADING----------
\begin{center}
  {\fontsize{26}{30}\selectfont\bfseries Jiaqi Xue}\\[5pt]
  \small
  \href{https://jqxue1999.github.io}{jqxue1999.github.io}
  \quad|\quad
  \href{mailto:jiaqi.xue@ucf.edu}{jiaqi.xue@ucf.edu}
  \quad|\quad
  \href{mailto:xuejiaqi001@gmail.com}{xuejiaqi001@gmail.com}
  \quad|\quad
  +1 (689) 837-6702
\end{center}

%----------BIOGRAPHY----------
\section{Biography}
Jiaqi Xue is a fourth-year Ph.D. candidate in Computer Science at the University of
Central Florida, advised by Prof.\ Qian Lou. He builds efficient and secure AI systems,
spanning efficient LLM serving and routing, secure and trustworthy AI, privacy-preserving
machine learning, and agentic code generation.

\vspace{4pt}
To date, he has published 17 papers in leading venues across machine learning
(NeurIPS, ICML, ICLR), NLP and vision (ACL, EMNLP, NAACL, ECCV), security and privacy
(IEEE S\&P, PoPETs, SaTML), and systems (DAC, PACT). He is the first author
of nine of them, including four at NeurIPS. He has interned
at Samsung Research America (SRA) and Amazon Web Services (AWS).

%----------EDUCATION----------
\section{Education}
\entry{University of Central Florida}{Orlando, FL}
      {Ph.D., Computer Science}{Jan.\ 2023 -- May 2027 (expected)}
\entry{Chongqing University}{Chongqing, China}
      {B.S., Computer Science}{Sep.\ 2018 -- Jun.\ 2022}

%----------EXPERIENCE----------
\section{Experience}
\entry{Amazon Web Services (AWS)}{Santa Clara, CA}
      {Applied Scientist Intern, Kiro Science}{May 2026 -- Aug.\ 2026}
\entrydesc{Multi-agent system (MAS) protocols and harnesses for code generation.}
\entry{Samsung Research America (SRA)}{Mountain View, CA}
      {Research Scientist Intern}{May 2024 -- Aug.\ 2024}
\entrydesc{Trustworthy LLMs.}
\entry{University of Central Florida (UCF)}{Orlando, FL}
      {Graduate Research Assistant}{Jan.\ 2023 -- Present}

%----------CONTRIBUTIONS----------
\section{Selected Research Contributions}
\begin{itemize}[leftmargin=1.1em,itemsep=1pt,topsep=0pt,parsep=0pt]
  \item \textbf{Efficient LLM serving and routing.} \textbf{R2-Router (ICML 2026)} lets
    a router reason about which model to use and how much computation to spend; it was
    the first method to push the RouterArena leaderboard beyond a score of 70.
    \textbf{Chain-of-Route (NeurIPS 2026)} extends routing from single queries to
    stateful, multi-turn agentic interactions, and \textbf{HW-Router (DAC 2026)} makes routing hardware-aware for scalable
    multi-LLM serving.
  \item \textbf{Secure and trustworthy AI.} \textbf{TrojLLM (NeurIPS 2023)} introduced the first
    automated prompt-injected Trojan attack on large language models. His other work
    covers backdoor attacks with \textbf{BadFair (EMNLP 2024)} and \textbf{TrojFSP (NAACL
    2024)}, defenses with \textbf{\mbox{SSL-Cleanse} (ECCV 2024)} and certified robustness
    with \textbf{CR-UTP (ACL 2024)}, and robust private inference with \textbf{RobPI
    (SaTML 2026)}. In LLM watermarking, \textbf{PRO (NeurIPS 2026)} enables
    precise and robust watermarks for open-source LLMs, and a \textbf{cross-lingual
    benchmark (EMNLP 2025)} evaluates the robustness of text watermarks.
  \item \textbf{Privacy-preserving machine learning.} \textbf{DictPFL (NeurIPS 2025)}
    makes federated learning on FHE-encrypted gradients efficient, and
    \textbf{CipherPrune (ICLR 2025)} scales hybrid MPC--FHE private Transformer
    inference. He also works on verifiable
    encrypted computation with \textbf{DataSeal (IEEE S\&P 2025)}, FHE acceleration with
    \textbf{\mbox{BoostCom} (PACT 2024)}, and FHE for general AI computation with a
    \textbf{systematization study (PoPETs 2026)}.
  \item \textbf{Agentic code generation.} \textbf{FHE-Coder (ICLR 2026)} enables secure
    agentic code generation for fully homomorphic encryption, a direction he extended
    to multi-agent protocols and harnesses at \textbf{AWS Kiro}.
\end{itemize}

%----------AWARDS----------
\section{Awards}
\begin{tabular}{@{}l@{\hspace{1.2em}}l@{}}
  2026 & DAC Young Fellow, Design Automation Conference.\\
  2025 & Faculty Cluster Initiative (FCI) Scholarship, University of Central Florida.\\
  2024 & Top Reviewer Award, NeurIPS.\\
  2023 & Scholar Award, NeurIPS.\\
\end{tabular}

%----------PUBLICATIONS----------
\section{Publications \texorpdfstring{{\normalfont\small\color{black}($^*$ equal contribution)}}{}}

\pub{20}{\me, Mengxin Zheng, Heng Huang, and Qian Lou,
  \ptitle{https://jqxue1999.github.io/files/Chain_of_Route__State_Aware_LLM_Routing_for_Multi_Turn_Conversations.pdf}{Chain-of-Route:
  State-Aware LLM Routing for Multi-Turn Conversations}. In Annual Conference on Neural
  Information Processing Systems (NeurIPS) 2026.}

\pub{19}{\me, Yifei Zhao, Mansour Al Ghanim, Shangqian Gao, Ruimin Sun, Qian Lou, and
  Mengxin Zheng, \ptitle{https://arxiv.org/pdf/2510.23891}{PRO: Enabling Precise and
  Robust Text Watermark for Open-Source LLMs}. In Annual Conference on Neural Information
  Processing Systems (NeurIPS) 2026.}

\pub{18}{\me, Qian Lou, Jiarong Xing, and Heng Huang,
  \ptitle{https://arxiv.org/pdf/2602.02823}{R2-Router: A New Paradigm for LLM Routing
  with Reasoning}. In International Conference on Machine Learning (ICML) 2026.}

\pub{17}{Mayank Kumar, \me, Mengxin Zheng, and Qian Lou,
  \ptitle{https://openreview.net/pdf?id=4F1py5vQXm}{FHE-Coder: Secure Agentic Code
  Generation for Fully Homomorphic Encryption}. In International Conference on Learning
  Representations (ICLR) 2026.}

\pub{16}{Ahasan Kabir, \me, Mengxin Zheng, and Qian Lou,
  \ptitle{https://arxiv.org/pdf/2608.14575}{HW-Router: Hardware-Aware Routing for
  Scalable Multi-LLM Serving}. In Design Automation Conference (DAC) 2026.}

\pub{15}{\me, Xin Xin, Wei Zhang, Mengxin Zheng, Qianqian Song, Minxuan Zhou, Yushun Dong,
  Dongjie Wang, Xun Chen, Jiafeng Xie, Liqiang Wang, David Mohaisen, Hongyi Wu, and Qian
  Lou, \ptitle{https://arxiv.org/pdf/2504.11604}{SoK: Can Fully Homomorphic Encryption
  Support General AI Computation? A Functional and Cost Analysis}. In Proceedings on
  Privacy Enhancing Technologies (PoPETs) 2026.}

\pub{14}{\me, Mengxin Zheng, and Qian Lou,
  \ptitle{https://arxiv.org/pdf/2602.19918}{RobPI: Robust Private Inference against
  Malicious Client}. In IEEE Conference on Secure and Trustworthy Machine Learning
  (SaTML) 2026.}

\pub{13}{\me, Mayank Kumar, Yuzhang Shang, Shangqian Gao, Rui Ning, Mengxin Zheng,
  Xiaoqian Jiang, and Qian Lou, \ptitle{https://arxiv.org/pdf/2510.21086}{DictPFL:
  Efficient and Private Federated Learning on Encrypted Gradients}. In Annual Conference on
  Neural Information Processing Systems (NeurIPS) 2025.}

\pub{12}{Mansour Al Ghanim, \me, Rochana Prih Hastuti, Mengxin Zheng, Yan Solihin, and
  Qian Lou, \ptitle{https://aclanthology.org/2025.findings-emnlp.390.pdf}{Evaluating the
  Robustness and Accuracy of Text Watermarking Under Real-World Cross-Lingual
  Manipulations}. In Conference on Empirical Methods in Natural Language
  Processing (EMNLP) 2025.}

\pub{11}{Yancheng Zhang, \me, Mengxin Zheng, Mimi Xie, Mingzhe Zhang, Lei Jiang, and Qian
  Lou, \ptitle{https://arxiv.org/pdf/2502.16782}{CipherPrune: Efficient and Scalable
  Private Transformer Inference}. In International Conference on Learning
  Representations (ICLR) 2025.}

\pub{10}{Muhammad Husni Santriaji, \me, Yancheng Zhang, Qian Lou, and Yan Solihin,
  \ptitle{https://www.computer.org/csdl/proceedings-article/sp/2025/223600a078/21B7RlxV5Ly}{DataSeal:
  Ensuring the Verifiability of Private Computation on Encrypted Data}. In IEEE
  Symposium on Security and Privacy (S\&P) 2025.}

\pub{9}{\me, Qian Lou, and Mengxin Zheng,
  \ptitle{https://aclanthology.org/2024.findings-emnlp.484.pdf}{BadFair: Backdoored
  Fairness Attacks with Group-conditioned Triggers}. In Conference on
  Empirical Methods in Natural Language Processing (EMNLP) 2024.}

\pub{8}{\me$^*$, Mengxin Zheng$^*$, Zihao Wang, Xun Chen, Qian Lou, Lei Jiang, and
  Xiaofeng Wang,
  \ptitle{https://www.ecva.net/papers/eccv_2024/papers_ECCV/papers/11933.pdf}{SSL-Cleanse:
  Trojan Detection and Mitigation in Self-Supervised Learning}. In European Conference
  on Computer Vision (ECCV) 2024.}

\pub{7}{Qian Lou, \me$^*$, Xin Liang$^*$, Yancheng Zhang, Rui Xie, and Mengxin Zheng,
  \ptitle{https://aclanthology.org/2024.findings-acl.588.pdf}{CR-UTP: Certified
  Robustness against Universal Text Perturbations on Large Language Models}. In Annual Meeting of the Association for Computational Linguistics (ACL) 2024.}

\pub{6}{Mengxin Zheng, \me, Xun Chen, Yanshan Wang, Qian Lou, and Lei Jiang,
  \ptitle{https://aclanthology.org/2024.naacl-long.64.pdf}{TrojFSP: Trojan Insertion in
  Few-shot Prompt Tuning}. In Annual Conference of the North American Chapter of the
  Association for Computational Linguistics (NAACL) 2024.}

\pub{5}{Ardhi Wiratama Baskara Yudha, \me, Qian Lou, Huiyang Zhou, and Yan Solihin,
  \ptitle{https://arxiv.org/pdf/2407.07308}{BoostCom: Towards Efficient Universal Fully
  Homomorphic Encryption by Boosting the Word-wise Comparisons}. In International
  Conference on Parallel Architectures and Compilation Techniques (PACT) 2024.}

\pub{4}{\me, Mengxin Zheng, Yi Sheng, Lei Yang, Qian Lou, and Lei Jiang,
  \ptitle{https://arxiv.org/pdf/2312.10508}{TrojFair: Trojan Fairness Attacks}. In ACM
  CCS Workshop on Large AI Systems and Models with Privacy and Safety Analysis (LAMPS)
  2024.}

\pub{3}{\me, Yancheng Zhang, Yanshan Wang, Xueqiang Wang, Hao Zheng, and Qian Lou,
  \ptitle{https://arxiv.org/pdf/2409.16675}{CryptoTrain: Fast Secure Training on
  Encrypted Dataset}. In ACM CCS Workshop on Large AI Systems and Models with Privacy
  and Safety Analysis (LAMPS) 2024.}

\pub{2}{\me, Mengxin Zheng, Ting Hua, Yilin Shen, Yepeng Liu, Ladislau B\"ol\"oni, and
  Qian Lou,
  \ptitle{https://proceedings.neurips.cc/paper_files/paper/2023/file/cf04d01a0e76f8b13095349d9caca033-Paper-Conference.pdf}{TrojLLM:
  A Black-box Trojan Prompt Attack on Large Language Models}. In Annual Conference on Neural
  Information Processing Systems (NeurIPS) 2023.}

\pub{1}{\me, Chentian Ma, Li Li, and Xuan Wen, Multiple EffNet/ResNet Architectures for
  Melanoma Classification. In International Conference on Computer Engineering and
  Application (ICCEA) 2021.}

%----------PREPRINTS----------
\section{Preprints}
\pub{P4}{\me, Yanjun Wang, Xiangci Li, Lingbo Mo, Aritra Sengupta, Shweta Garg, Murali
  Krishna Ramanathan, and Myeongsoo Kim, AECP: Artifact-Exclusive Communication Protocol
  for Multi-Agent Code Generation. Under review.}

\pub{P3}{\me, Mengxin Zheng, and Qian Lou, From Capability to Alignment: Repurposing MoE
  Routing for Safety. Under review.}

\pub{P2}{\me, Yifei Zhao, Mengxin Zheng, Xun Chen, Fan Yao, Yan Solihin, and Qian Lou,
  \ptitle{https://arxiv.org/pdf/2507.03278}{Securing Transformer-based AI Execution via
  Unified TEE and Crypto-protected Accelerators}. Under review.}

\pub{P1}{\me, Mengxin Zheng, Yebowen Hu, Fei Liu, Xun Chen, and Qian Lou,
  \ptitle{https://arxiv.org/pdf/2406.00083}{BadRAG: Identifying Vulnerabilities in
  Retrieval Augmented Generation of Large Language Models}. Under review.}

%----------TEACHING----------
\section{Teaching}
Graduate Teaching Assistant, University of Central Florida:\par\vspace{2pt}
\begin{tabular}{@{}l@{\hspace{1.2em}}l@{}}
  CDA5106 & Advanced Computer Architecture (Fall 2023, Fall 2024, Fall 2025)\\
  CAP6614 & Current Topics in Machine Learning (Spring 2024, Spring 2025)\\
  CIS3360 & Security in Computing (Spring 2026)\\
  CDA3103 & Computer Logic and Organization (Summer 2023)\\
\end{tabular}

%----------SERVICE----------
\section{Service}
Reviewer: NeurIPS, ICML, ICLR, AISTATS, AAAI, IJCAI, ACL, EMNLP, CVPR, ICCV, TMLR.

\vspace{12pt}
{\footnotesize\color{rulegray} Last updated: September 2026}

\end{document}
